\documentclass{article}
\usepackage[margin=2cm]{geometry}
\usepackage{amsmath}
\usepackage{graphicx}
\usepackage{booktabs}
\usepackage[utf8]{inputenc}
\usepackage{ragged2e}
\setlength{\emergencystretch}{3em}
\begin{document}
\sloppy

\noindent\hfill\includegraphics[height=1.8em]{assets/figure_header_logo.png}\par\noindent\hrule\vspace{4pt}

among never-takers with D = 1 (as defiers do not exist). Likewise, 0 corresponds to the direct effect on always-takers with D = O. Indeed, the results in Table
suggest that both parameters are consistently estimated with the change-in-changes
model (Panel A.).


\subsection*{Proof of Theorem 1}


1.1Average direct effect under d = 1 conditional on D = 1
andM(1)= 0


() = () - (0) = ()0  1d    
0] = E[Yi - Qoo(Yo)|D = 1, M = 0]. Using the observational rule, we obtain
E[Yi(1,0)|D = 1, M(1) = 0] = E[Yi|D = 1,M = 0]. Accordingly, we have to
 y [0 =  = l)0] =[ =() = )]  Mos
proof. Denote the inverse of h(d, m,t,u) by h-1(d, m,t;y), which exists because of
the strict monotonicity required in Assumption 1.Under Assumptions 1 and 3a,the
conditional potential outcome distribution function equals


\begin{equation}
\begin{array} { r l } { F _ { Y _ { t } ( d, 0 ) | D = 1, M = 0 } ( y ) \stackrel { A! } { = } \operatorname* { P r } ( h ( d, m, t, U ) \leq y | D = 1, M = 0, T = t ), } & { } \\ { = } & { \mathrm { P r } ( U \leq h ^ { - 1 } ( d, m, t ; y ) | D = 1, M = 0, T = t ), } & { } \\ & { \qquad = F _ { U | 1 0 } ( d,}\end{array}
\end{equation}


for d,d' E \{O,1\}. We use these quantities in the following. First, evaluating
FY(0.0)D=1,M=0(y) at h(0, 0, 1,u) gives





\begin{equation}
h ( 0, 0, 1, u ) = F _ { Y _ { 1 } ( 0, 0 ) | D = 1, M = 0 } ^ { - 1 } ( F _ { U | 1 0 } ( u ) ). \qquad ( 1. 2 )
\end{equation}

\end{document}
